% Algebra 2 Lesson 2-6 fast-track packet: The Quadratic Formula (enVision A2, Example 1 and Example 3 only)
% Quiz map (2-6 Lesson Quiz, a2_2-6_TE.tex): items 1-2 -> Example 1 (A real, B non-real); item 3 -> Example 3;
% item 4 -> Example 4 (cut; rehearsed by the Practice block, TE practice 30-31); item 5 -> Example 5 (cut; discriminant = 0 from Example 3).
% Shared preamble for the "fast" (hacked-down) Algebra 2 lesson packets, Oct 2026.
% Each packet: % Shared preamble for the "fast" (hacked-down) Algebra 2 lesson packets, Oct 2026.
% Each packet: % Shared preamble for the "fast" (hacked-down) Algebra 2 lesson packets, Oct 2026.
% Each packet: \input{fast-packet-template} then \begin{document} ... \end{document}.
\documentclass[12pt]{article}
\usepackage[margin=0.7in]{geometry}
\usepackage{amsmath,amssymb}
\usepackage{booktabs,tabularx,array}
\usepackage{enumitem}
\usepackage[T1]{fontenc}
\usepackage[utf8]{inputenc}
\usepackage{lmodern}
\usepackage{tikz}
\usepackage{pgfplots}
\pgfplotsset{compat=1.18}
\usepackage{xcolor}
\setlist{nosep,leftmargin=1.4em}
\setlength{\parindent}{0pt}
\setlength{\parskip}{6pt}
\renewcommand{\arraystretch}{1.35}
\newcolumntype{Y}{>{\raggedright\arraybackslash}X}
\pagestyle{empty}

% blanks and workspace
\newcommand{\blank}[1][2.6cm]{\rule[-2pt]{#1}{0.5pt}}
\newcommand{\lines}[1][1]{\par\vspace{2pt}\foreach \i in {1,...,#1}{\rule{0pt}{22pt}\rule{\linewidth}{0.4pt}\par}}
\newcommand{\wspart}[2]{\vspace{4pt}{\large\bfseries #1: #2}\par\vspace{-6pt}\rule{\textwidth}{0.4pt}\par}

% header: \packethead{Lesson 1-3}{Piecewise-Defined Functions}
\newcommand{\packethead}[2]{%
{\LARGE\bfseries #2}\hfill Name \blank[5cm]\\[2pt]
Algebra 2 \quad #1 (fast track) \hfill Date \blank[2.5cm] \quad Period \blank[1.2cm]\par\vspace{2pt}}

% boxed goals: \goals{learning goal}{language goal}
\newcommand{\goals}[2]{%
\fbox{\parbox{\dimexpr\textwidth-2\fboxsep-2\fboxrule}{%
\textbf{Learning goal:} #1\\[2pt]
\textbf{Language goal:} #2}}\par}

% quiz target box: \quiztarget{text}
\newcommand{\quiztarget}[1]{%
\vspace{4pt}\colorbox{gray!12}{\parbox{\dimexpr\textwidth-2\fboxsep}{\textbf{On the quiz:} #1}}\par}

% IXL table row helper: code & skill & which quiz question it feeds
\newcommand{\ixlrow}[3]{\texttt{#1} & #2 & #3 \\}

% one grid style for every graph; override xmin/xmax/... per plot
\pgfplotsset{sheetgrid/.style={
  axis lines=middle, grid=both, grid style={gray!45},
  xmin=-10.5, xmax=10.5, ymin=-10.5, ymax=10.5,
  xtick={-10,-8,...,10}, ytick={-10,-8,...,10},
  x=0.28cm, y=0.28cm,
  ticklabel style={font=\scriptsize, fill=white, inner sep=1pt},
  xlabel={$x$}, ylabel={$y$},
  axis line style={thick}, clip=true}}
% unit grid for small windows
\pgfplotsset{unitgrid/.style={sheetgrid,
  xmin=-5.5, xmax=5.5, ymin=-5.5, ymax=5.5,
  xtick={-5,...,5}, ytick={-5,...,5}, x=0.5cm, y=0.5cm}}
% open and closed endpoints
\tikzset{closed/.style={circle, fill, inner sep=1.6pt},
         open/.style={circle, draw, fill=white, inner sep=1.4pt}}
 then \begin{document} ... \end{document}.
\documentclass[12pt]{article}
\usepackage[margin=0.7in]{geometry}
\usepackage{amsmath,amssymb}
\usepackage{booktabs,tabularx,array}
\usepackage{enumitem}
\usepackage[T1]{fontenc}
\usepackage[utf8]{inputenc}
\usepackage{lmodern}
\usepackage{tikz}
\usepackage{pgfplots}
\pgfplotsset{compat=1.18}
\usepackage{xcolor}
\setlist{nosep,leftmargin=1.4em}
\setlength{\parindent}{0pt}
\setlength{\parskip}{6pt}
\renewcommand{\arraystretch}{1.35}
\newcolumntype{Y}{>{\raggedright\arraybackslash}X}
\pagestyle{empty}

% blanks and workspace
\newcommand{\blank}[1][2.6cm]{\rule[-2pt]{#1}{0.5pt}}
\newcommand{\lines}[1][1]{\par\vspace{2pt}\foreach \i in {1,...,#1}{\rule{0pt}{22pt}\rule{\linewidth}{0.4pt}\par}}
\newcommand{\wspart}[2]{\vspace{4pt}{\large\bfseries #1: #2}\par\vspace{-6pt}\rule{\textwidth}{0.4pt}\par}

% header: \packethead{Lesson 1-3}{Piecewise-Defined Functions}
\newcommand{\packethead}[2]{%
{\LARGE\bfseries #2}\hfill Name \blank[5cm]\\[2pt]
Algebra 2 \quad #1 (fast track) \hfill Date \blank[2.5cm] \quad Period \blank[1.2cm]\par\vspace{2pt}}

% boxed goals: \goals{learning goal}{language goal}
\newcommand{\goals}[2]{%
\fbox{\parbox{\dimexpr\textwidth-2\fboxsep-2\fboxrule}{%
\textbf{Learning goal:} #1\\[2pt]
\textbf{Language goal:} #2}}\par}

% quiz target box: \quiztarget{text}
\newcommand{\quiztarget}[1]{%
\vspace{4pt}\colorbox{gray!12}{\parbox{\dimexpr\textwidth-2\fboxsep}{\textbf{On the quiz:} #1}}\par}

% IXL table row helper: code & skill & which quiz question it feeds
\newcommand{\ixlrow}[3]{\texttt{#1} & #2 & #3 \\}

% one grid style for every graph; override xmin/xmax/... per plot
\pgfplotsset{sheetgrid/.style={
  axis lines=middle, grid=both, grid style={gray!45},
  xmin=-10.5, xmax=10.5, ymin=-10.5, ymax=10.5,
  xtick={-10,-8,...,10}, ytick={-10,-8,...,10},
  x=0.28cm, y=0.28cm,
  ticklabel style={font=\scriptsize, fill=white, inner sep=1pt},
  xlabel={$x$}, ylabel={$y$},
  axis line style={thick}, clip=true}}
% unit grid for small windows
\pgfplotsset{unitgrid/.style={sheetgrid,
  xmin=-5.5, xmax=5.5, ymin=-5.5, ymax=5.5,
  xtick={-5,...,5}, ytick={-5,...,5}, x=0.5cm, y=0.5cm}}
% open and closed endpoints
\tikzset{closed/.style={circle, fill, inner sep=1.6pt},
         open/.style={circle, draw, fill=white, inner sep=1.4pt}}
 then \begin{document} ... \end{document}.
\documentclass[12pt]{article}
\usepackage[margin=0.7in]{geometry}
\usepackage{amsmath,amssymb}
\usepackage{booktabs,tabularx,array}
\usepackage{enumitem}
\usepackage[T1]{fontenc}
\usepackage[utf8]{inputenc}
\usepackage{lmodern}
\usepackage{tikz}
\usepackage{pgfplots}
\pgfplotsset{compat=1.18}
\usepackage{xcolor}
\setlist{nosep,leftmargin=1.4em}
\setlength{\parindent}{0pt}
\setlength{\parskip}{6pt}
\renewcommand{\arraystretch}{1.35}
\newcolumntype{Y}{>{\raggedright\arraybackslash}X}
\pagestyle{empty}

% blanks and workspace
\newcommand{\blank}[1][2.6cm]{\rule[-2pt]{#1}{0.5pt}}
\newcommand{\lines}[1][1]{\par\vspace{2pt}\foreach \i in {1,...,#1}{\rule{0pt}{22pt}\rule{\linewidth}{0.4pt}\par}}
\newcommand{\wspart}[2]{\vspace{4pt}{\large\bfseries #1: #2}\par\vspace{-6pt}\rule{\textwidth}{0.4pt}\par}

% header: \packethead{Lesson 1-3}{Piecewise-Defined Functions}
\newcommand{\packethead}[2]{%
{\LARGE\bfseries #2}\hfill Name \blank[5cm]\\[2pt]
Algebra 2 \quad #1 (fast track) \hfill Date \blank[2.5cm] \quad Period \blank[1.2cm]\par\vspace{2pt}}

% boxed goals: \goals{learning goal}{language goal}
\newcommand{\goals}[2]{%
\fbox{\parbox{\dimexpr\textwidth-2\fboxsep-2\fboxrule}{%
\textbf{Learning goal:} #1\\[2pt]
\textbf{Language goal:} #2}}\par}

% quiz target box: \quiztarget{text}
\newcommand{\quiztarget}[1]{%
\vspace{4pt}\colorbox{gray!12}{\parbox{\dimexpr\textwidth-2\fboxsep}{\textbf{On the quiz:} #1}}\par}

% IXL table row helper: code & skill & which quiz question it feeds
\newcommand{\ixlrow}[3]{\texttt{#1} & #2 & #3 \\}

% one grid style for every graph; override xmin/xmax/... per plot
\pgfplotsset{sheetgrid/.style={
  axis lines=middle, grid=both, grid style={gray!45},
  xmin=-10.5, xmax=10.5, ymin=-10.5, ymax=10.5,
  xtick={-10,-8,...,10}, ytick={-10,-8,...,10},
  x=0.28cm, y=0.28cm,
  ticklabel style={font=\scriptsize, fill=white, inner sep=1pt},
  xlabel={$x$}, ylabel={$y$},
  axis line style={thick}, clip=true}}
% unit grid for small windows
\pgfplotsset{unitgrid/.style={sheetgrid,
  xmin=-5.5, xmax=5.5, ymin=-5.5, ymax=5.5,
  xtick={-5,...,5}, ytick={-5,...,5}, x=0.5cm, y=0.5cm}}
% open and closed endpoints
\tikzset{closed/.style={circle, fill, inner sep=1.6pt},
         open/.style={circle, draw, fill=white, inner sep=1.4pt}}

\begin{document}

\packethead{Lesson 2-6}{The Quadratic Formula}
\goals{I can solve quadratic equations using the Quadratic Formula and use the discriminant to predict the number and type of solutions.}{I can explain how many solutions a quadratic equation has and what type they are, using the words discriminant, radicand, real, and non-real.}

\quiztarget{(1) Solve with the Quadratic Formula; pick the simplified solutions (Example 1A). (2) Bonus: the same with a negative radicand, so the solutions have $i$ (Example 1B). (3) Use the discriminant to tell how many roots and what type (Example 3). (4) Will a tossed ball reach a height? (full Savvas quiz only) (5) Which $b$ make the discriminant $0$? (full Savvas quiz only)}

\fbox{\parbox{\dimexpr\textwidth-2\fboxsep-2\fboxrule}{%
\textbf{The Quadratic Formula.} For any quadratic equation $ax^2+bx+c=0$ with $a\neq0$:
\vspace{-4pt}\[x=\frac{-b\pm\sqrt{b^2-4ac}}{2a}\]\vspace{-10pt}

The expression under the radical, $b^2-4ac$, is the \textbf{radicand}. It is also called the \textbf{discriminant}.}}

\wspart{Example 1}{Solve Quadratic Equations}
What are the solutions to the equation? Write the equation as $ax^2+bx+c=0$, name $a$, $b$, $c$, then substitute. Use parentheses around negative numbers.

\begin{minipage}[t]{0.48\textwidth}
\textbf{A.} $3x^2-4x-9=0$

$a=3$, \ $b=-4$, \ $c=-9$
\begin{align*}
x&=\frac{-(-4)\pm\sqrt{(-4)^2-4(3)(-9)}}{2(3)}\\
&=\frac{4\pm\sqrt{16+108}}{6}=\frac{4\pm\sqrt{124}}{6}\\
&=\frac{4\pm2\sqrt{31}}{6} \qquad (\sqrt{124}=\sqrt{4\cdot31})\\
&=\frac{2\pm\sqrt{31}}{3}
\end{align*}
The radicand is $124>0$: \textbf{two real} solutions,
\[x=\frac{2+\sqrt{31}}{3}\quad\text{and}\quad x=\frac{2-\sqrt{31}}{3}.\]
\end{minipage}\hfill
\begin{minipage}[t]{0.48\textwidth}
\textbf{B.} $x^2-9x+27=0$

$a=1$, \ $b=-9$, \ $c=27$
\begin{align*}
x&=\frac{-(-9)\pm\sqrt{(-9)^2-4(1)(27)}}{2(1)}\\
&=\frac{9\pm\sqrt{81-108}}{2}=\frac{9\pm\sqrt{-27}}{2}\\
&=\frac{9\pm i\sqrt{27}}{2} \qquad (\sqrt{-1}=i)\\
&=\frac{9\pm3i\sqrt{3}}{2} \qquad (\sqrt{27}=3\sqrt3)
\end{align*}
The radicand is $-27<0$: \textbf{two non-real} solutions,
\[x=\frac{9+3i\sqrt{3}}{2}\quad\text{and}\quad x=\frac{9-3i\sqrt{3}}{2}.\]
\end{minipage}

\textbf{Say it:} In Part B the radicand is \blank[1.6cm]. It is \blank[2.2cm] (positive / negative), so the solutions are \blank[2.2cm] (real / non-real) because \dots\lines[2]

\newpage

\wspart{Try It 1}{Solve using the Quadratic Formula}
\begin{minipage}[t]{0.48\textwidth}
\textbf{a.} $2x^2+6x+3=0$

\vspace{2pt}
$a=\blank[1cm]$ \quad $b=\blank[1cm]$ \quad $c=\blank[1cm]$

\vspace{2pt}
$x=\dfrac{-(\phantom{00})\pm\sqrt{(\phantom{00})^2-4(\phantom{00})(\phantom{00})}}{2(\phantom{00})}$
\lines[8]
\vspace{4pt}
Radicand: \blank[1.6cm]\\[4pt]
Real or non-real? \blank[2.2cm]

$x=\blank[5cm]$
\end{minipage}\hfill
\begin{minipage}[t]{0.48\textwidth}
\textbf{b.} $3x^2-2x+7=0$

\vspace{2pt}
$a=\blank[1cm]$ \quad $b=\blank[1cm]$ \quad $c=\blank[1cm]$

\vspace{2pt}
$x=\dfrac{-(\phantom{00})\pm\sqrt{(\phantom{00})^2-4(\phantom{00})(\phantom{00})}}{2(\phantom{00})}$
\lines[8]
\vspace{4pt}
Radicand: \blank[1.6cm]\\[4pt]
Real or non-real? \blank[2.2cm]

$x=\blank[5cm]$
\end{minipage}

In part b, where does the $i$ come from? Use the word \emph{radicand}.\lines[2]

\vspace{4pt}
\fbox{\parbox{\dimexpr\textwidth-2\fboxsep-2\fboxrule}{%
\textbf{Challenge (DOK 3, self-paced).} Seaira wants to use the Quadratic Formula to solve the equation $x^4+5x^2-5=0$. Is this possible? If so, describe the steps she should follow.
\lines[3]}}

\newpage

\wspart{Example 3}{Identify the Number of Real-Number Solutions}
How can you determine the number and type of roots for a quadratic equation? Graph each equation. Then use the Quadratic Formula to find the roots.

\vspace{2pt}
\noindent
\begin{tabular}{@{}p{0.315\textwidth}p{0.315\textwidth}p{0.315\textwidth}@{}}
\centering $y=2x^2-7x+3$ & \centering $y=4x^2+12x+9$ & \centering\arraybackslash $y=x^2+2x+8$ \\
\centering
\begin{tikzpicture}
% TE: x-window [-5,5], y-window [-6,6], unit grid, ticks -4,-2,2,4 on both axes, O at origin, arrowed curve
\begin{axis}[unitgrid, xmin=-5.5, xmax=5.5, ymin=-6.5, ymax=6.5,
  xtick={-4,-2,2,4}, ytick={-4,-2,2,4}, x=0.45cm, y=0.36cm,
  minor xtick={-5,...,5}, minor ytick={-6,...,6}, minor grid style={gray!45},
  yticklabel style={anchor=east, xshift=-3pt}]
\addplot[very thick, blue, domain=-0.3:3.8, samples=80, <->] {2*x^2-7*x+3};
\node[closed] at (axis cs:0.5,0) {};
\node[closed] at (axis cs:3,0) {};
\node[closed] at (axis cs:0,3) {};
\node[closed] at (axis cs:1.75,-3.125) {};
\node[font=\scriptsize, inner sep=1pt, anchor=north east] at (axis cs:0,0) {$O$};
\node[font=\scriptsize, fill=white, inner sep=0.5pt, anchor=north west] at (axis cs:0.5,-3.5) {$(1.75,-3.125)$};
\end{axis}
\end{tikzpicture}
&
\centering
\begin{tikzpicture}
% TE: x-window [-5,5], y-window [-5,7], unit grid, x labels -4,-2,O,2,4, y labels -4,-2,2,4,6, arrowed curve
% y tick labels sit right of the y-axis so the curve (all left of it) does not cross them
\begin{axis}[unitgrid, xmin=-5.5, xmax=5.5, ymin=-5.5, ymax=7.5,
  xtick={-4,-2,2,4}, ytick={-4,-2,2,4,6}, x=0.45cm, y=0.36cm,
  minor xtick={-5,...,5}, minor ytick={-5,...,7}, minor grid style={gray!45},
  yticklabel style={anchor=west, xshift=4pt}]
\addplot[very thick, blue, domain=-2.8:-0.2, samples=80, <->] {4*x^2+12*x+9};
\node[closed] at (axis cs:-1.5,0) {};
\node[font=\scriptsize, inner sep=1pt, anchor=north west] at (axis cs:0,0) {$O$};
\node[font=\scriptsize, fill=white, inner sep=0.5pt, anchor=north] at (axis cs:-1.5,-1.35) {$(-1.5,0)$};
\end{axis}
\end{tikzpicture}
&
\centering\arraybackslash
\begin{tikzpicture}
% TE: x-window [-5,5], y-window [-2,10], unit grid, x labels -4,-2,O,2,4, y labels 2,4,6,8, arrowed curve
\begin{axis}[unitgrid, xmin=-5.5, xmax=5.5, ymin=-2.5, ymax=10.5,
  xtick={-4,-2,2,4}, ytick={2,4,6,8}, x=0.45cm, y=0.36cm,
  minor xtick={-5,...,5}, minor ytick={-2,...,10}, minor grid style={gray!45},
  yticklabel style={anchor=west, xshift=4pt}]
\addplot[very thick, blue, domain=-2.4:0.4, samples=80, <->] {x^2+2*x+8};
\node[closed] at (axis cs:-1,7) {};
\node[closed] at (axis cs:0,8) {};
\node[font=\scriptsize, inner sep=1pt, anchor=north east] at (axis cs:0,0) {$O$};
\node[font=\scriptsize, fill=white, inner sep=0.5pt, anchor=north east] at (axis cs:-1.1,6.7) {$(-1,7)$};
\end{axis}
\end{tikzpicture}
\\[2pt]
$2x^2-7x+3=0$ & $4x^2+12x+9=0$ & $x^2+2x+8=0$ \\[2pt]
$x=\dfrac{7\pm\sqrt{(-7)^2-4(2)(3)}}{2(2)}$ & $x=\dfrac{-12\pm\sqrt{(12)^2-4(4)(9)}}{2(4)}$ & $x=\dfrac{-2\pm\sqrt{(2)^2-4(1)(8)}}{2(1)}$ \\[8pt]
$x=\dfrac{7\pm\sqrt{25}}{4}=\dfrac{7\pm5}{4}$ & $x=\dfrac{-12\pm\sqrt{0}}{8}=\dfrac{-12}{8}$ & $x=\dfrac{-2\pm\sqrt{-28}}{2}=\dfrac{-2\pm2i\sqrt7}{2}$ \\[8pt]
$x=3$ \ and \ $x=\tfrac12$ & $x=-\tfrac32$ & $x=-1\pm i\sqrt{7}$ \\[4pt]
\small The radicand is $25$. Adding and subtracting $\sqrt{25}$ gives two different roots. The graph shows \textbf{two} $x$-intercepts: \textbf{two real} roots. &
\small The radicand is $0$. Adding and subtracting $\sqrt0$ gives one root. The graph shows \textbf{one} $x$-intercept (the vertex): \textbf{one real} root. &
\small The radicand is $-28$. Its square root is imaginary. The graph shows \textbf{no} $x$-intercepts: \textbf{two non-real} roots. \\
\end{tabular}

\vspace{4pt}
\fbox{\parbox{\dimexpr\textwidth-2\fboxsep-2\fboxrule}{%
The \textbf{discriminant} is the radicand, $b^2-4ac$ ($a$, $b$, $c$ real, $a\neq0$). It tells the nature of the roots:\\[2pt]
\hspace*{1em}$b^2-4ac>0$: two real roots \quad $b^2-4ac=0$: one real root \quad $b^2-4ac<0$: two non-real roots}}

\textbf{Say it:} For $x^2+2x+8=0$ the discriminant is \blank[1.4cm]. It is \blank[2cm], so the equation has \blank[1.2cm] \blank[2cm] roots, and the graph \blank[3.5cm] the $x$-axis.

\wspart{Try It 3}{Describe the nature of the solutions for each equation}
\begin{minipage}[t]{0.48\textwidth}
\textbf{a.} $16x^2+8x+1=0$

$a=\blank[0.9cm]$ \ $b=\blank[0.9cm]$ \ $c=\blank[0.9cm]$

$b^2-4ac=$\lines[2]
The discriminant is \blank[1.6cm] ($>0$, $=0$, $<0$), so there is/are \blank[3.2cm] solution(s).
\end{minipage}\hfill
\begin{minipage}[t]{0.48\textwidth}
\textbf{b.} $2x^2-5x+6=0$

$a=\blank[0.9cm]$ \ $b=\blank[0.9cm]$ \ $c=\blank[0.9cm]$

$b^2-4ac=$\lines[2]
The discriminant is \blank[1.6cm] ($>0$, $=0$, $<0$), so there is/are \blank[3.2cm] solution(s).
\end{minipage}

\newpage

% R5: parallels quiz item 4 (Richard's ball, h(t) = -5t^2 + 10t + 6, reach 12 m?). TE practice items 30-31 (SEE EXAMPLE 4), verbatim.
\wspart{Practice}{Will the ball reach that height?}
\begin{minipage}[t]{0.58\textwidth}
\vspace{0pt}
Deon throws a ball into the air. The height, $h$, of the ball, in meters, at time $t$ seconds is modeled by the function $h(t)=-5t^2+t+4$.

\textbf{a.} When will the ball hit the ground? (Hit the ground: $h=0$.)\lines[3]

\textbf{b.} Will the ball reach a height of 5 meters?

Set $h(t)=5$: \blank[4cm]

Standard form: \blank[4cm] $=0$

$a=\blank[0.9cm]$ \ $b=\blank[0.9cm]$ \ $c=\blank[0.9cm]$ \quad $b^2-4ac=\blank[1.8cm]$

So: \lines[1]
\end{minipage}\hfill
\begin{minipage}[t]{0.38\textwidth}
\vspace{0pt}
\begin{center}
\begin{tikzpicture}
% TE practice 30: x-window [-3,3], y-window [-2,5], unit grid, labels x=-2,O,2 and y=2,4, curve ends arrowed downward
% (TE labels the axes x and y; the packet uses t and h to match h(t))
\begin{axis}[unitgrid, xmin=-3.5, xmax=3.5, ymin=-2.5, ymax=5.5,
  xtick={-2,2}, ytick={2,4}, x=0.85cm, y=0.6cm,
  minor xtick={-3,...,3}, minor ytick={-2,...,5}, minor grid style={gray!45},
  xlabel={$t$}, ylabel={$h$}]
\addplot[very thick, blue, domain=-1:1.2, samples=80, <->] {-5*x^2+x+4};
\node[font=\scriptsize, inner sep=1pt, anchor=north east] at (axis cs:0,0) {$O$};
\end{axis}
\end{tikzpicture}
\end{center}
\end{minipage}

\wspart{IXL}{Practice for the quiz}
Do the starred skills first. \textbf{SmartScore 70 or above on each starred ($\bigstar$) skill replaces your quiz score.}

\vspace{2pt}
\begin{tabularx}{\textwidth}{l >{\hsize=.8\hsize}Y >{\hsize=1.2\hsize}Y}
\toprule
\textbf{Code} & \textbf{Skill} & \textbf{Feeds} \\
\midrule
\ixlrow{YQH $\bigstar$}{Solve a quadratic equation using the quadratic formula}{Quiz 1; Quiz 2 (bonus). Example 1, Try It 1}
\ixlrow{QHK $\bigstar$}{Using the discriminant}{Quiz 3. Example 3, Try It 3 (also Quiz 4, 5: full Savvas quiz only)}
\ixlrow{5VV}{Introduction to complex numbers}{Optional warm-up, no quiz score: square roots of negatives (Example 1B)}
\ixlrow{CEU}{Solve quadratic equations: word problems}{Optional: Practice page; Quiz 4 (full Savvas quiz only)}
\bottomrule
\end{tabularx}

\vspace{2pt}
Prerequisites if you need review (no quiz score): \texttt{CJC} solve by factoring, \texttt{NPH} completing the square.

\wspart{Before the quiz}{Finish these sentences}
To use the Quadratic Formula, first write the equation in the form \blank[3.5cm] $=0$.

The radicand of the Quadratic Formula, $b^2-4ac$, is called the \blank[3cm].

If the discriminant is positive, the equation has \blank[2.8cm]. If it is zero, it has \blank[2.8cm].

If the discriminant is negative, the equation has two \blank[2.5cm] solutions, and the graph does not \blank[2.5cm] the $x$-axis.

\end{document}
